% Generated by `taskgen course examples` from factoring.congruence_of_squares@trivial (seed 1).
% Do not edit: change examples.yaml and regenerate.
\begin{example}[комбинирование соотношений, тривиальная и удачная зависимости]\label{ex:squares-combination}
Для $N = 3649$ найдены соотношения $a_i^2 \equiv c_i \pmod{N}$, в которых числа $c_i$ раскладываются по факторной базе $\{2, 3, 5, 7, 11\}$:\[\begin{array}{c|r|r|l}i & a_i & c_i & \\ \hline 1 & 79 & 2592 & = 2^{5} \cdot 3^{4} \\ 2 & 125 & 1029 & = 3 \cdot 7^{3} \\ 3 & 121 & 45 & = 3^{2} \cdot 5 \\ 4 & 86 & 98 & = 2 \cdot 7^{2} \\ 5 & 83 & 3240 & = 2^{3} \cdot 3^{4} \cdot 5 \\ 6 & 67 & 840 & = 2^{3} \cdot 3 \cdot 5 \cdot 7 \\ 7 & 73 & 1680 & = 2^{4} \cdot 3 \cdot 5 \cdot 7 \\ \end{array}\]С помощью исключения Гаусса над $\mathbb{F}_2$ найти наборы соотношений, произведение $c_i$ в которых --- полный квадрат, и для каждого (по порядку) вычислить $X = \prod a_i \bmod N$, $Y = \sqrt{\prod c_i} \bmod N$ и $\gcd(X - Y, N)$, пока не будет найден нетривиальный делитель $N$.

\par\noindent\textit{Решение.} Векторы показателей по модулю $2$ (столбцы --- простые $2, 3, 5, 7, 11$):\[\begin{array}{c|ccccc} & 2 & 3 & 5 & 7 & 11 \\ \hline v_{1} & 1 & 0 & 0 & 0 & 0 \\ v_{2} & 0 & 1 & 0 & 1 & 0 \\ v_{3} & 0 & 0 & 1 & 0 & 0 \\ v_{4} & 1 & 0 & 0 & 0 & 0 \\ v_{5} & 1 & 0 & 1 & 0 & 0 \\ v_{6} & 1 & 1 & 1 & 1 & 0 \\ v_{7} & 0 & 1 & 1 & 1 & 0 \\ \end{array}\]Исключение по столбцам слева направо (ведущая строка --- первая ещё не использованная строка с единицей в столбце; её прибавляем к остальным неиспользованным строкам с единицей в этом столбце, запоминая, из каких $v_i$ составлена строка):\begin{itemize}\item столбец $2$: ведущая строка $1$; к строкам $4, 5, 6$ прибавляем её;\item столбец $3$: ведущая строка $2$; к строкам $6, 7$ прибавляем её;\item столбец $5$: ведущая строка $3$; к строкам $5, 6, 7$ прибавляем её;\item столбец $7$: единиц в неиспользованных строках нет;\item столбец $11$: единиц в неиспользованных строках нет;\end{itemize}Не ставшие ведущими строки обнулились; они дают зависимости $v_{1} + v_{4} = 0$, $v_{1} + v_{3} + v_{5} = 0$, $v_{1} + v_{2} + v_{3} + v_{6} = 0$, $v_{2} + v_{3} + v_{7} = 0$.\par\noindent \textbf{Набор $\{1, 4\}$.} $\prod c_i = 2^{6} \cdot 3^{4} \cdot 7^{2}$, $X = 79 \cdot 86 \bmod 3649 = 3145$, $Y = 2^{3} \cdot 3^{2} \cdot 7 \bmod 3649 = 504$; $\gcd(X - Y, N) = \gcd(2641, 3649) = 1$ --- тривиальный делитель: $X \equiv -Y \pmod{N}$ ($X + Y = 3649$), переходим к следующей зависимости.\par\noindent \textbf{Набор $\{1, 3, 5\}$.} $\prod c_i = 2^{8} \cdot 3^{10} \cdot 5^{2}$, $X = 79 \cdot 121 \cdot 83 \bmod 3649 = 1564$, $Y = 2^{4} \cdot 3^{5} \cdot 5 \bmod 3649 = 1195$; $\gcd(X - Y, N) = \gcd(369, 3649) = 41$: $X^2 \equiv Y^2$, но $X \not\equiv \pm Y \pmod{N}$, делитель найден.\[N = 41 \cdot 89.\]
\end{example}
